\section{Related literature}
\label{sec:literature}

DEFINE-UK sits at the junction of three literatures: stock-flow consistent macroeconomic modelling, ecological macroeconomics, and the climate-and-financial-stability agenda. A fourth --- the practice of replication itself --- is what this paper contributes to.

\subsection{Stock-flow consistent modelling}

The stock-flow consistent (SFC) tradition descends from \citet{godley2007}, whose organising insight is that accounting comes before theory: every monetary flow leaves one sector's account and enters another's, and every flow accumulates into a stock, so that a transactions matrix sums to zero in every row and a balance-sheet matrix sums to zero in every instrument row. Nothing appears from nowhere and nothing leaks out. The discipline is unusually strong for a macroeconomic modelling framework, because it converts a large class of specification errors into arithmetic failures that a matrix will not close over.

Two surveys map the field the model belongs to. \citet{caverzasi2015} trace the post-Keynesian SFC programme from its Cambridge origins through the empirical and open-economy extensions, and note the characteristic methodological trade: SFC models buy accounting completeness at the cost of large parameter sets calibrated rather than estimated, which makes their published tables --- initial values, parameter provenance --- an unusually large part of what a reader must be able to check. \citet{nikiforos2017} survey the same literature with emphasis on how such models are closed, and on the recurring result that a demand-led closure delivers larger fiscal multipliers than a supply-constrained one, because output is not returned to a potential path by assumption. Section~\ref{sec:calibration} shows that this is not an abstract point for DEFINE-UK: it is one of the three reasons its green-investment multiplier cannot be read against short-horizon published figures.

The empirical relevance of these surveys to the present paper is direct. In an SFC model calibrated rather than estimated, the published parameter and initial-value tables \emph{are} the specification. Section~\ref{sec:reimplementation} is in large part a report on what happens when one takes that seriously enough to transcribe those tables in full and demand that the printed equations reproduce the printed values.

\subsection{Ecological macroeconomics and the DEFINE framework}

The specific model replicated here belongs to the DEFINE family. \citet{dafermos2017} introduced the stock-flow-fund framework that DEFINE implements: a synthesis of Godley--Lavoie monetary accounting with the material flow-fund accounting of ecological economics, so that the same model tracks money and matter under separate but coupled consistency conditions. \citet{dafermos2018} extended it to the climate--finance channel, showing how physical damages and transition dynamics feed through firm leverage and bank balance sheets to credit availability, and how monetary policy interacts with that loop --- the mechanism DEFINE-UK's credit-rationing and default blocks inherit.

DEFINE-UK \citep{georgedafermos2026manual} is the United Kingdom calibration and extension of that framework. Its distinguishing content relative to the generic case is the pair of UK-specific structural blocks described in Section~\ref{sec:model}: a housing block that tracks the dwelling stock by energy efficiency, so that retrofit policy operates through the stock rather than through an assumed elasticity, and a power-generation sector with separate fossil and non-fossil capital, an investment trigger tied to capacity utilisation, and --- new in version~1.1 --- regulatory instruments that directly depreciate fossil capital.

\subsection{Climate policy, stranded assets and financial stability}

The reason to want balance sheets in a climate model is the transition-risk agenda that \citet{carney2015} named: the losses from a rapid, credible decarbonisation fall not on the emitting capital alone but on the claims written against it, held by institutions whose solvency and lending behaviour then feed back into the real economy. \citet{battiston2017} made that concrete with a network stress test showing how climate-policy shocks propagate from exposed sectors through equity and credit holdings to the wider financial system. DEFINE-UK's version~1.1 regulation instruments are an in-model implementation of that mechanism: a regulatory scenario writes down power-sector capital, and the resulting change in leverage, illiquidity and credit rationing is a modelled consequence rather than an add-on. Whether it does so \emph{correctly} is not something this paper can establish for the scenarios, for reasons Section~\ref{sec:validation} sets out.

\subsection{Replication as a research output}

Finally, this paper belongs to the replication literature, in the sense \citet{dewald1986} established: that published results should be reproducible from published materials, and that the systematic attempt to do so is itself a finding. Their project's central lesson --- that the binding constraint is usually the completeness of what was published, not the willingness of the replicator --- describes this exercise closely. Here the constraint is threefold and unusually explicit: the code is unlicensed, the scenario results are unpublished, and the manual's parameter tables and equations do not everywhere agree with each other. Each is recorded as a ceiling with a stated route to lifting it, rather than routed around.

Within the PolicyEngine Macro suite the same standard has been applied to a central bank's structural VAR \citep{ahmadi2026boesvar} and to the OBR's macroeconometric model \citep{ahmadi2026obr}. The methodological point those papers share with this one is the separation, at the point of claim, between checks that can fail and checks that cannot. A tolerance chosen after seeing the deviation it admits, an identity that holds by construction, a band fitted to a verbal quantity --- these are worth reporting, but they are not evidence of agreement, and a replication that presents them as such is overstating its result. Section~\ref{sec:validation} labels each one where it appears.
